\documentclass{article}
\usepackage[utf8]{inputenc}
\usepackage{amsmath}
\usepackage{amsfonts}
\usepackage{amssymb}
\usepackage{geometry}
\usepackage{tcolorbox}
\usepackage{hyperref}
\geometry{a4paper, margin=1in}

\title{Module 08: Production Pipelines \& Collaborative Git Workflows - Part 1}
\author{Cybernauts 2026 Datathon Campaign}
\date{May 2026}

\begin{document}

\maketitle

\begin{abstract}
Transitioning from chaotic Jupyter notebooks to structured, modular script architectures is the hallmark of professional data engineering. In high-pressure datathons, unorganized notebooks cause cell-execution errors, variable overwrites, and tracking failures. This note defines the software engineering principles required to modularize an experimental machine learning workflow into a robust, single-responsibility script directory under a master execution runner.
\end{abstract}

\tableofcontents
\newpage

\section{The Failure of Notebook-Centric Workflows}
In the early phases of data exploration, Jupyter notebooks are excellent for quick visualizations and ad-hoc feature testing. However, relying on them for your final production pipeline introduces severe architectural weaknesses.

\subsection{The Hidden State Trap}
Notebooks maintain an out-of-order execution state memory. If you execute Cell 15, scroll up to modify Cell 4, and then execute Cell 20, your workspace variables depend entirely on your manual clicking history. This lack of linear execution leads to silent data bugs and makes reproducing cross-validation scores nearly impossible.

\subsection{Collaborative Bottlenecks}
When working in a hackathon team, merging raw notebook JSON files via Git is an engineering nightmare. Version controlling a file that updates metadata every time a cell is executed leads to constant merge conflicts, slowing down collaborative iterations.

\section{The Single-Responsibility Directory Blueprint}
To build an elite pipeline, you must migrate your codebase into a standardized, single-responsibility directory layout. By separating concerns, team members can simultaneously develop features, tune models, and alter paths without blocking each other.

The production configuration utilizes an isolated source engine directory (\texttt{src/}) controlled by a root master file:

\begin{verbatim}
.
├── data/
│   ├── train.csv
│   └── test.csv
├── src/
│   ├── __init__.py
│   ├── config.py
│   ├── preprocessing.py
│   ├── features.py
│   └── models.py
├── main.py
└── submission.csv
\end{verbatim}

\subsection{\texttt{src/config.py}: The Control Panel}
This configuration script houses all global constants, absolute and relative file paths, algorithmic seed configurations, and hyperparameter mapping spaces.
\begin{itemize}
    \item \textbf{Engineering Rule:} No raw numbers or file paths should be hardcoded anywhere else in your codebase. If you need to switch from a 5-fold to a 10-fold split, or redirect the data ingestion path to a new file variant (e.g., \texttt{TRAIN\_PATH = "data/train\_v2.csv"}), you alter exactly one file parameter.
\end{itemize}

\subsection{\texttt{src/preprocessing.py}: The Sanitization Ingestion Engine}
This module contains pure python functions tasked exclusively with handling raw data loading, missing value imputation, and basic categorical conversions. It handles structural formatting without engineering advanced synthetic properties.

\subsection{\texttt{src/features.py}: The Mathematical Synthesis Engine}
The dedicated factory for advanced feature engineering. This module ingests clean data frames and runs sequential functional operations to output high-value mathematical combinations, target encodings, out-of-fold metrics, and group-level context aggregations.

\subsection{\texttt{src/models.py}: The Cross-Validation and Ensembling Layer}
This script encapsulates your multi-fold training framework. It coordinates fold splitting, instantiates tree estimators, tracks and displays localized out-of-fold evaluation metrics, and handles your ensembling, blending, and threshold calibration logic.

\section{The Master Executor Architecture}
The orchestration of your modular software components is handled exclusively by a unified entry point script located at the project root directory: \texttt{main.py}.

\subsection{Linear Ingestion Pipelines}
The role of \texttt{main.py} is to serve as an explicit, step-by-step pipeline runner. It imports functions from your \texttt{src/} scripts and executes them sequentially in a clean, reproducible stream:

\begin{equation}
    \text{Raw Data} \xrightarrow{\text{preprocessing.py}} \text{Sanitized Data} \xrightarrow{\text{features.py}} \text{Engineered Array} \xrightarrow{\text{models.py}} \text{Submission Out}
\end{equation}

Running the simple command \texttt{python main.py} triggers this end-to-end execution path completely from scratch. It loads files, handles column imputations, derives interaction values, runs cross-validated gradient boosting models, and safely exports an optimized, calibrated final output submission file.

\newpage
\section{Practice Problems}

\textbf{P1:} Design a sample implementation of a \texttt{src/config.py} script that follows the single-responsibility principle, defining configurations for data paths, a random state seed, a 5-fold cross-validation split parameter, and a placeholder baseline dictionary for LightGBM hyperparameters.
\\
\textit{Answer:}
\begin{verbatim}
# src/config.py
import os

# Structural File Vectors
DATA_DIR = "data"
TRAIN_PATH = os.path.join(DATA_DIR, "train.csv")
TEST_PATH = os.path.join(DATA_DIR, "test.csv")
SUBMISSION_OUTPUT = "submission.csv"

# Global Reproducibility Parameters
RANDOM_SEED = 2026
NUM_FOLDS = 5

# Estimator Parameters
LGBM_PARAMS = {
    'objective': 'binary',
    'metric': 'auc',
    'learning_rate': 0.05,
    'num_leaves': 31,
    'random_state': RANDOM_SEED,
    'n_estimators': 1000
}
\end{verbatim}

\newline
\textbf{P2:} During a pipeline refactoring sprint, a teammate writes a function inside \texttt{src/features.py} that reads the training data file directly from a hardcoded string path, rather than importing it. Explain why this approach violates modular engineering principles and identify which script should handle data path definitions.
\\
\textit{Answer: This approach violates the single-responsibility principle by creating tight coupling and a hardcoding anti-pattern. If the data paths are scattered across processing and feature scripts, updating a data source location requires searching through every file to refactor the paths, which can introduce silent bugs. Data paths must live exclusively inside \texttt{src/config.py}, and functions within \texttt{src/features.py} should only accept in-memory data objects (such as pandas DataFrames) passed directly from the main orchestration loop.}

\newline
\textbf{P3:} Draw out a concise conceptual schematic layout of the \texttt{main.py} execution routine using Python code structure, showcasing how it orchestrates data passing between individual single-responsibility files.
\\
\textit{Answer:}
\begin{verbatim}
# main.py
import src.config as config
from src.preprocessing import load_and_clean_data
from src.features import generate_synthetic_features
from src.models import run_cross_validation_ensemble

def main():
    print("[1/4] Extracting and Cleaning Raw Inputs...")
    df_train, df_test = load_and_clean_data(config.TRAIN_PATH, config.TEST_PATH)

    print("[2/4] Executing Feature Engineering Ingestion...")
    X_train, X_test, y_train = generate_synthetic_features(df_train, df_test)

    print("[3/4] Running Cross-Validation Ensembles & Calibration...")
    oof_predictions, test_predictions = run_cross_validation_ensemble(
        X_train, X_test, y_train, config.NUM_FOLDS, config.LGBM_PARAMS
    )

    print("[4/4] Generating Submission Output Track...")
    test_predictions.to_csv(config.SUBMISSION_OUTPUT, index=False)
    print("Pipeline executed successfully. Output secured.")

if __name__ == "__main__":
    main()
\end{verbatim}

\end{document}
